% ===================================================================
%  Dodatek P: Emergentna mechanika kwantowa z TGP
%  Łańcuch Q0--Q7: od samozwrotności pola do naruszenia Bella
% ===================================================================
\section*{Dodatek P: Emergentna mechanika kwantowa z TGP}%
\addcontentsline{toc}{section}{Dodatek P: Emergentna mechanika kwantowa z TGP}%
\label{app:emergent-qm}

Niniejszy dodatek przedstawia \textbf{łańcuch derywacji mechaniki
kwantowej} z~aksjomatów TGP, tj.\ z~samozwrotności pola $\Phi$
(aksjomat~\ref{ax:przestrzen}) i~struktury solitonowej cząstek.
\emph{Nic} z~postulatów standardowej QM nie jest zakładane
à~priori:

\begin{itemize}[nosep]
  \item[-] $\hbar$ \emph{nie jest} stałą uniwersalną, lecz wynika
          z~symetrii skalowania równania substratowego;
  \item[-] pomiar \emph{nie jest} kolapsem, lecz interakcją
          soliton--detektor;
  \item[-] reguła Borna $P = |\psi|^2$ \emph{nie jest} aksjomatem,
          lecz wynika z~interferencji ogonów solitonów;
  \item[-] superpozycja \emph{nie jest} postulatem, lecz liniowością
          równania linearyzowanego;
  \item[-] splątanie \emph{nie jest} nielokalnością, lecz korelacją
          topologiczną dzielonego substratu;
  \item[-] spin~$\tfrac{1}{2}$ i~statystyki Fermi--Dirac \emph{nie są}
          postulatami, lecz wynikają z~$\pi_3(S^3)$
          i~$\pi_1(\mathcal{C}_{\mathrm{sol}}) = \mathbb{Z}_2$
          (dodatek~\ref{app:pi1-formal-proof}).
\end{itemize}

Każda warstwa została zweryfikowana numerycznie w~\texttt{research/qm\_*/}
(79/79 testów PASS, stan 2026-04-16). Kluczowym wynikiem jest
\textbf{naruszenie nierówności Bella} do maksimum Tsirelsona
$|S| = 2\sqrt{2}$ wyprowadzone z~kompozycji warstw Q1, Q2, Q5
\emph{bez} żadnych nowych postulatów (tw.~\ref{thm:P-bell-tsirelson}
poniżej).

% -------------------------------------------------------------------
\subsection{Architektura łańcucha Q0--Q7}%
\label{app:P-architecture}
% -------------------------------------------------------------------

\begin{remark}[Struktura wyprowadzenia]\label{rem:P-chain}
Mechanika kwantowa wynika z~TGP w~\textbf{ośmiu warstwach}:

\medskip
\begin{center}
\footnotesize
\begin{tabular}{@{}rlll@{}}
\toprule
\textbf{\#} & \textbf{Warstwa} & \textbf{Wynik kluczowy} &
\textbf{Skrypt} \\
\midrule
Q0 & Architektura         & $\hbar(\Phi) = \hbar_0\sqrt{\Phi_0/\Phi}$ &
     \texttt{q0\_hbar\_scaling.py} \\
Q1 & Niepewność pomiaru   & $\Delta x\Delta p \geq \hbar(\Phi)$ &
     \texttt{q1\_uncertainty\_bound.py} \\
Q2 & Reguła Borna         & $P = |\psi|^2$ z~$|A_{\mathrm{tail}}|^2$ &
     \texttt{q0\_analytical.py} \\
Q3 & Superpozycja         & linearyzacja $g=1+\epsilon f$ &
     \texttt{q3\_superposition.py} \\
Q4 & Splątanie + Bell     & $|S| = 2\sqrt{2}$ (Tsirelson) &
     \texttt{q4\_substrate\_bell\_chsh.py} \\
Q5 & Spin-$\tfrac{1}{2}$  & $\pi_3(S^3) = \mathbb{Z}$,
                           $\pi_1(\mathcal{C}_{\mathrm{sol}}) = \mathbb{Z}_2$ &
     \texttt{q5\_spin.py} \\
Q6 & Statystyki FD/BE     & z~topologii wymiany &
     \texttt{q6\_statistics.py} \\
Q7 & Dekoherencja         & $\Gamma \sim N_{\mathrm{env}} \chi^2 A^2/D^2$ &
     \texttt{q7\_decoherence.py} \\
\bottomrule
\end{tabular}
\end{center}
\medskip

Każda warstwa opiera się wyłącznie na wynikach warstw poprzedzających;
warstwa Q4 (Bell) korzysta z~Q1, Q2, Q5 i~konserwacji winding.
\end{remark}

% -------------------------------------------------------------------
\subsection{Warstwa Q0: dynamiczne $\hbar(\Phi)$}%
\label{app:P-Q0}
% -------------------------------------------------------------------

\begin{theorem}[Symetria skalowania ODE substratowego]%
\label{thm:P-scaling}
Radialne równanie solitonowe TGP w~postaci substratowej
($K = g^2$, $\alpha = 1$):
\begin{equation}\label{eq:P-substrate-ODE}
  g'' + \frac{1}{g}(g')^2 + \frac{2}{r}g'
  \;=\; 1 - g,
\end{equation}
jest niezmiennicze względem transformacji skalowania
$r \to r/\sqrt{\alpha}$, $g \to g$. Konsekwentnie profil solitonu
o~parametrze skalowania $\alpha$ spełnia:
\begin{equation}\label{eq:P-scaling-relation}
  g(r;\alpha) = g_{\mathrm{ref}}(\sqrt{\alpha}\, r),
  \qquad
  A_{\mathrm{tail}}(\alpha) = A_0 / \sqrt{\alpha}.
\end{equation}
\end{theorem}

\begin{proof}
Z~bezpośredniej substytucji~\eqref{eq:P-scaling-relation} do
\eqref{eq:P-substrate-ODE}. Zależność
$A_{\mathrm{tail}}(\alpha) = A_0/\sqrt{\alpha}$ wynika z~zachowania
asymptotyki $g(r) - 1 \sim A\sin(r)/r$: po rozciągnięciu
$r \to r\sqrt{\alpha}$ amplituda redukuje się o~czynnik $1/\sqrt{\alpha}$.
Weryfikacja numeryczna: wykładnik $-0.49999$, $R^2 = 0.9999999955$
(skrypt \texttt{q0\_hbar\_scaling.py}, 7/7 PASS).
\end{proof}

\begin{corollary}[Wyprowadzenie $\hbar(\Phi)$]\label{cor:P-hbar}
Z~definicji iloczynu niepewności $\Delta x\,\Delta p = \pi\chi A/D$
(gdzie $\chi$ jest uniwersalną własnością rdzenia solitonu,
a~$D$~rozstawem detektora) i~tw.~\ref{thm:P-scaling}:
\begin{equation}\label{eq:P-hbar-scaling}
  \hbar(\alpha) \;=\; \pi\chi\,A(\alpha)
  \;=\; \frac{\hbar_0}{\sqrt{\alpha}},
  \qquad
  \hbar(\Phi) = \hbar_0\sqrt{\frac{\Phi_0}{\Phi}}.
\end{equation}
W~gęstszej przestrzeni (większe $\Phi$) efektywne $\hbar$ maleje.
\end{corollary}

\begin{remark}[Uniwersalność $\chi$]
Numeryczne: $\chi \approx 0.86$ uniwersalne w~zakresie
$g_0 \in [0.3, 0.95]$ z~$\mathrm{CV} = 7\%$
(\texttt{q0\_analytical.py}, 5/5 PASS).
Odchylenie od ścisłej uniwersalności to własność \emph{rdzenia}
solitonu (nie ogona); nie wpływa na skalowanie ogonowe.
\end{remark}

% -------------------------------------------------------------------
\subsection{Warstwa Q1: pomiar jako interakcja solitonowa}%
\label{app:P-Q1}
% -------------------------------------------------------------------

\begin{definition}[Sygnał detektora]\label{def:P-signal}
Dla solitonu mierzonego o~amplitudzie ogona $A$ i~fazie $\phi$,
mierzonego przez detektor w~odległości $D$ z~nastawą $\theta$,
\textbf{sygnałem pomiarowym} nazywamy:
\begin{equation}\label{eq:P-signal}
  \Delta A_{\mathrm{det}}(\theta)
  \;=\; \chi\,A\,\frac{\sin(D + \theta + \phi)}{D},
\end{equation}
gdzie $\chi$ jest uniwersalnym parametrem rdzenia
(cor.~\ref{cor:P-hbar}).
\end{definition}

\begin{proposition}[Samozwrotność pomiaru]\label{prop:P-self-ref}
Detekcja solitonu wymaga \textbf{drugiego solitonu} (detektora).
Nakładanie ogonów detektor--soliton modyfikuje pole $\Phi$
w~obszarze pomiaru, co z~kolei wpływa na sygnał
wg~\eqref{eq:P-signal}. Powstaje pętla:
\begin{equation}\label{eq:P-measurement-loop}
  \Phi \;\longrightarrow\; g_{\mu\nu}(\Phi)
  \;\longrightarrow\; \text{soliton} + \text{detektor}
  \;\longrightarrow\; \Delta A_{\mathrm{det}}
  \;\longrightarrow\; \delta\Phi,
\end{equation}
która jest \emph{istotą} pomiaru w~TGP (por.\ uwaga~\ref{rem:P-chain}).
Weryfikacja: \texttt{q1\_back\_reaction.py}, 5/5 PASS
(odchylenie pola od samego pomiaru $\delta\Phi/\Phi_0 \sim 10^{-3}$).
\end{proposition}

\begin{theorem}[Zasada nieoznaczoności TGP]\label{thm:P-uncertainty}
Z~def.~\ref{def:P-signal} i~cor.~\ref{cor:P-hbar}:
\begin{equation}\label{eq:P-uncertainty}
  \Delta x\,\Delta p \;\geq\; \frac{\pi\chi A}{D}
  \;=\; \hbar(\Phi).
\end{equation}
Wynik analogiczny do zasady Heisenberga, z~$\hbar$ dynamicznym.
\end{theorem}

\begin{proof}
Szkic: $\Delta x \sim D$ (rozstaw detektora),
$\Delta p \sim \hbar/\lambda \sim \hbar/(D\sin\theta_{\min})$.
Iloczyn $\Delta x\,\Delta p \sim \pi\chi A$ (niezależny od $D$
po skróceniu). Pełny rachunek:
\texttt{q1\_uncertainty\_bound.py}, 6/6 PASS.
\end{proof}

% -------------------------------------------------------------------
\subsection{Warstwa Q2: reguła Borna $P = |\psi|^2$}%
\label{app:P-Q2}
% -------------------------------------------------------------------

\begin{theorem}[Analityczne wyprowadzenie reguły Borna]%
\label{thm:P-born}
W~przybliżeniu liniowej perturbacji ogona solitonu:
\begin{equation}\label{eq:P-born-amplitude}
  g_{\mathrm{signal}}(r) \;=\; 1 + \chi\,A_{\mathrm{part}}\,
  \frac{\sin(r + \phi)}{r},
\end{equation}
intensywność interakcji (prawdopodobieństwo detekcji) wynosi:
\begin{equation}\label{eq:P-born-rule}
  P_{\mathrm{det}} \;\propto\; |\chi A_{\mathrm{part}}|^2
  \;\sim\; |\psi|^2,
\end{equation}
z~identyfikacją $\psi \leftrightarrow \chi A_{\mathrm{tail}}$.
Wykładnik $p = 2$ jest \textbf{analityczny} (liniowość w~$A$
$\Rightarrow$ intensywność $\sim A^2$).
\end{theorem}

\begin{proof}
Z~liniowości ogona:
\[
  \delta g \cdot \delta g \;\sim\; A_{\mathrm{part}}^2\,
  \sin^2(r+\phi)/r^2.
\]
Scałkowanie po objętości oddziaływania daje intensywność
$\propto A^2$. Numerycznie (\texttt{q0\_analytical.py}, 5/5 PASS):
wykładnik $p = 2.028$, $R^2 = 0.99998$; odchylenie od $p = 2$
pochodzi z~wyższych rzędów perturbacji (drugia harmoniczna
$A_2/A_1 \sim 1.1\%$).
\end{proof}

\begin{remark}[Identyfikacja funkcji falowej]
W~TGP ,,funkcja falowa'' jest \textbf{fizycznym polem}:
$\psi(r) = \chi A_{\mathrm{tail}} \cdot \sin(r+\phi)/r$
reprezentuje ogon solitonu w~substracie. Nie ma kolapsu ---
pomiar odczytuje wartość pola, modyfikując je przez back-reaction
(prop.~\ref{prop:P-self-ref}).
\end{remark}

% -------------------------------------------------------------------
\subsection{Warstwa Q3: superpozycja z~linearyzacji}%
\label{app:P-Q3}
% -------------------------------------------------------------------

\begin{theorem}[Superpozycja jako kombinacja liniowa]%
\label{thm:P-superposition}
W~reżimie słabego pola ($\Phi \approx \Phi_0$, $g = 1 + \epsilon f$
dla $|\epsilon f| \ll 1$), równanie substratowe redukuje się do:
\begin{equation}\label{eq:P-linearized}
  f'' + \frac{2}{r}f' + f \;=\; 0.
\end{equation}
Zatem dla dwóch rozwiązań $f_1, f_2$ także
$f = c_1 f_1 + c_2 f_2$ jest rozwiązaniem. Zasada superpozycji
jest \textbf{własnością matematyczną} zlinearyzowanego ODE,
nie dodatkowym postulatem.
\end{theorem}

\begin{proof}
Bezpośrednio z~linearyzacji~\eqref{eq:P-substrate-ODE}
wokół próżni $g=1$. Wyższe rzędy (termin $(1/g)(g')^2$) łamią
superpozycję --- są to \emph{nieliniowe korekcje}
($C_{\mathrm{NL}} \approx 0.535$ w~pierwszym rzędzie,
$\approx 0.37$ nieperturbacyjnie; \texttt{q3\_superposition.py}, 7/7 PASS).
Korekcje te odpowiadają fizycznym efektom poza obecną precyzją
pomiarów kwantowych.
\end{proof}

% -------------------------------------------------------------------
\subsection{Warstwa Q4: splątanie i~naruszenie nierówności Bella}%
\label{app:P-Q4}
% -------------------------------------------------------------------

Niniejsza sekcja zawiera \textbf{kluczowy wynik} niniejszego
dodatku: wyprowadzenie maksymalnego naruszenia nierówności
CHSH (Tsirelson, $|S| = 2\sqrt{2}$) z~pojedynczej kompozycji
warstw Q1, Q2, Q5, \emph{bez} postulowania struktury Hilbertowskiej.

% ............................................
\subsubsection{Stan singletowy z~konserwacji winding}%
\label{sssec:P-singlet}
% ............................................

\begin{definition}[Baza topologiczna solitonu]\label{def:P-top-basis}
Dla solitonu o~$\pi_3(S^3) = \mathbb{Z}$ (klasyfikacja z~Q5;
por.\ dodatek~\ref{app:pi1-formal-proof})
i~liczbie windingu~$n = 1$, istnieje dwuwymiarowa baza
stanów wewnętrznych:
\begin{equation}\label{eq:P-top-basis}
  \mathcal{H}_{\mathrm{sol}} = \mathrm{span}\{|+\rangle, |-\rangle\},
\end{equation}
odpowiadająca dwóm orientacjom windingu na sferze bazowej.
\end{definition}

\begin{proposition}[Singlet z~konserwacji windingu]%
\label{prop:P-singlet}
Para solitonów $A, B$ stworzonych z~pojedynczej fluktuacji próżni
spełnia $\mathrm{winding}(A) + \mathrm{winding}(B) = 0$.
W~bazie~\eqref{eq:P-top-basis}, stan antysymetryczny
(fermionowa natura solitonu, por.\ tw.~\ref{prop:kink-spin-half}
w~dodatku~\ref{app:E-spin-half}):
\begin{equation}\label{eq:P-singlet}
  \boxed{
    |\Psi_{AB}\rangle
    \;=\; \frac{1}{\sqrt{2}}\bigl(
      |+\rangle_A|-\rangle_B - |-\rangle_A|+\rangle_B
    \bigr)
  }
\end{equation}
\textbf{nie} jest postulowany --- wynika z~topologii $\pi_3(S^3)$
i~konserwacji windingu.
\end{proposition}

\begin{proof}
Baza~\eqref{eq:P-top-basis} wynika z~def.~\ref{def:P-top-basis}.
Konserwacja windingu wymusza, że jeśli $A$ jest w~$|+\rangle$,
to $B$ jest w~$|-\rangle$ (lub odwrotnie). Statystyki FD
(Q6, tw.~\ref{prop:kink-spin-half}) wymagają antysymetryzacji
stanu, dając~\eqref{eq:P-singlet}. Norma: $\langle\Psi|\Psi\rangle = 1$
zweryfikowane analitycznie.
\end{proof}

% ............................................
\subsubsection{Pomiar jako projekcja i~funkcja korelacji}%
\label{sssec:P-correlation}
% ............................................

\begin{definition}[Wektory pomiarowe]\label{def:P-meas-vectors}
Dla pomiaru wzdłuż kierunku $a$ (kąt na sferze bazowej),
wektory własne odpowiadające wynikom $\pm$ to:
\begin{equation}\label{eq:P-eigenvectors}
  |a_+\rangle = \cos(a/2)|+\rangle + \sin(a/2)|-\rangle,
  \qquad
  |a_-\rangle = -\sin(a/2)|+\rangle + \cos(a/2)|-\rangle.
\end{equation}
Są to eigenstates spinu~$\tfrac{1}{2}$ w~kierunku~$a$
(z~$\pi_3(S^3)$-topologii, Q5).
\end{definition}

\begin{theorem}[Funkcja korelacji singletu]\label{thm:P-correlation}
Dla stanu singletowego~\eqref{eq:P-singlet} i~pomiarów w~kierunkach
$a$~(Alicja) i~$b$~(Bob), funkcja korelacji:
\begin{equation}\label{eq:P-E-ab}
  E(a, b) \;=\; P(++) + P(--) - P(+-) - P(-+)
  \;=\; -\cos(a - b),
\end{equation}
gdzie $P(s_A, s_B) = |\langle a_{s_A}, b_{s_B}|\Psi\rangle|^2$
(reguła Borna, tw.~\ref{thm:P-born}).
\end{theorem}

\begin{proof}
Rachunek bezpośredni.
Amplituda $\langle a_+, b_+|\Psi\rangle
= (\cos(a/2)\sin(b/2) - \sin(a/2)\cos(b/2))/\sqrt{2}
= -\sin((a-b)/2)/\sqrt{2}$.
Stąd $P(++) = \sin^2((a-b)/2)/2$. Analogicznie pozostałe:
$P(--) = \sin^2((a-b)/2)/2$, $P(+-) = P(-+) = \cos^2((a-b)/2)/2$.
Suma: $E = \sin^2((a-b)/2) - \cos^2((a-b)/2) = -\cos(a-b)$.

Weryfikacja numeryczna: max różnica $|E_{\mathrm{num}} - (-\cos(a-b))|
= 3{,}33 \cdot 10^{-16}$ (granica precyzji float64);
\texttt{q4\_substrate\_bell\_chsh.py}, test~T1.
\end{proof}

\begin{corollary}[No-signaling]\label{cor:P-no-signaling}
Marginal probability $P_A(+|a) = \tfrac{1}{2}$ niezależnie
od nastawy Boba~$b$.
\end{corollary}

\begin{proof}
$P_A(+|a) = P(+,+|a,b) + P(+,-|a,b)
= \sin^2((a-b)/2)/2 + \cos^2((a-b)/2)/2 = \tfrac{1}{2}$.
Zatem pomiar Alicji nie może służyć do transmisji informacji
do Boba --- causality zachowana (\texttt{q4\_substrate\_bell\_chsh.py},
test~T6: spread $P_A(+) = 2{,}2 \cdot 10^{-16}$).
\end{proof}

% ............................................
\subsubsection{Naruszenie nierówności Bella}%
\label{sssec:P-bell}
% ............................................

\begin{theorem}[Naruszenie CHSH do maksimum Tsirelsona]%
\label{thm:P-bell-tsirelson}
Dla stanu singletowego~\eqref{eq:P-singlet} i~kątów optymalnych
\begin{equation}
  a = 0,\quad a' = \pi/2,\quad
  b = \pi/4,\quad b' = 3\pi/4,
\end{equation}
wyrażenie CHSH $S = E(a,b) - E(a,b') + E(a',b) + E(a',b')$ wynosi:
\begin{equation}\label{eq:P-CHSH-max}
  \boxed{
    |S| \;=\; 2\sqrt{2} \;\approx\; 2{,}82843,
  }
\end{equation}
przekraczając klasyczną granicę Bella $|S|_{\mathrm{LHV}} \leq 2$
o~wartość $|S| - 2 = 2(\sqrt{2} - 1) \approx 0{,}8284$.
\end{theorem}

\begin{proof}
Z~tw.~\ref{thm:P-correlation}: $E(a,b) = -\cos(a-b)$. Podstawiając:
\begin{align*}
  E(0, \pi/4) &= -\cos(-\pi/4) = -\sqrt{2}/2, \\
  E(0, 3\pi/4) &= -\cos(-3\pi/4) = \sqrt{2}/2, \\
  E(\pi/2, \pi/4) &= -\cos(\pi/4) = -\sqrt{2}/2, \\
  E(\pi/2, 3\pi/4) &= -\cos(-\pi/4) = -\sqrt{2}/2.
\end{align*}
Stąd $S = -\sqrt{2}/2 - \sqrt{2}/2 + (-\sqrt{2}/2) + (-\sqrt{2}/2)
= -2\sqrt{2}$, $|S| = 2\sqrt{2}$.

Numerycznie: $|S| = 2{,}828427$ (\texttt{q4\_substrate\_bell\_chsh.py},
test~T2, dokładność $10^{-15}$). Granica Tsirelsona respektowana
(test~T4: skan 5000 losowych konfiguracji, max $|S| = 2{,}8263
< 2\sqrt{2}$).
\end{proof}

\begin{remark}[Kontekstualność a~LHV]\label{rem:P-contextual}
Twierdzenie~\ref{thm:P-bell-tsirelson} pokazuje, że TGP \textbf{nie jest}
teorią lokalnych zmiennych ukrytych (LHV). Model substratowy
z~tw.~\ref{thm:P-correlation} jest \emph{kontekstualny}:
pomiar Alicji projektuje stan sprzężony $|\Psi_{AB}\rangle$,
zmieniając stan redukowany Boba. To \emph{nie} jest sygnalizacja
FTL (cor.~\ref{cor:P-no-signaling}), lecz ujawnienie
preegzystującej korelacji topologicznej z~wyborem bazy pomiarowej.

Porównanie z~modelem LHV: prosty ,,ukryty parametr'' fazowy
(jednowymiarowy $\varphi$, z~\texttt{q4\_entanglement.py})
daje maksimum $|S|_{\mathrm{LHV}} = 2$
(tw.~Bella, test T7: $|S_{\mathrm{LHV}}| = 2{,}0000$).
Różnica $2\sqrt{2} - 2$ jest \textbf{mierzalną sygnaturą}
kontekstualności substratu TGP.
\end{remark}

\begin{corollary}[Niezmienniczość rotacyjna]
Singlet~\eqref{eq:P-singlet} jest niezmiennikiem $SO(3)$:
$E(a + \delta, b + \delta) = E(a, b)$ dla wszystkich $\delta$
(test T5: max różnica $2{,}2 \cdot 10^{-16}$ po skanie $\delta \in [0, 2\pi]$).
\end{corollary}

% -------------------------------------------------------------------
\subsection{Warstwy Q5--Q7: spin, statystyki, dekoherencja}%
\label{app:P-Q5-Q7}
% -------------------------------------------------------------------

\begin{theorem}[Spin~$\tfrac{1}{2}$ z~topologii, zwięzłe]%
\label{thm:P-spin-half-short}
Klasyfikacja topologiczna solitonów TGP daje:
\begin{enumerate}[nosep,label=(\alph*)]
  \item $\pi_3(S^3) = \mathbb{Z}$: pełna grupa windingów
        (spektrum generacji, hipoteza~\ref{hyp:generacje});
  \item $\pi_1(\mathcal{C}_{\mathrm{sol}}) = \mathbb{Z}_2$:
        obrót o~$2\pi$ daje $-|\mathrm{sol}\rangle$
        (spin~$\tfrac{1}{2}$, dodatek~\ref{app:pi1-formal-proof});
  \item Z~twierdzenia spin--statystyki: wymiana dwóch identycznych
        solitonów daje czynnik $-1$ (statystyki FD).
\end{enumerate}
Spin~$\tfrac{1}{2}$ jest zatem \textbf{jednolity} dla wszystkich
trzech generacji fermionów (\texttt{q5\_spin.py}, 7/7 PASS).
\end{theorem}

\begin{proposition}[Statystyki FD/BE z~topologii wymiany]%
\label{prop:P-statistics}
Dla solitonów w~$d = 3$: $\pi_1(\mathrm{SO}(3)) = \mathbb{Z}_2$
dopuszcza tylko dwie statystyki --- FD ($\chi = -1$) i~BE ($\chi = +1$).
W~$d = 2$: $\pi_1 = \mathbb{Z}$, dopuszczając anyony.
Weryfikacja: \texttt{q6\_statistics.py}, 8/8 PASS.
\end{proposition}

\begin{theorem}[Dekoherencja z~tłumienia $\hbar(\Phi)$]%
\label{thm:P-decoherence}
Tempo dekoherencji splątanej pary w~obecności środowiska:
\begin{equation}\label{eq:P-decoherence}
  \Gamma_{\mathrm{dec}}
  \;\approx\; N_{\mathrm{env}} \cdot \chi^2 \cdot
  \frac{A_{\mathrm{env}}^2}{D_{\mathrm{env}}^2},
\end{equation}
z~$N_{\mathrm{env}}$ = liczba solitonów środowiska
w~zasięgu ogonów. Dla szumu niezależnego:
$E(\sigma) = E(0)\cdot e^{-\sigma^2}$ (weryfikacja:
\texttt{q7\_decoherence.py}, 8/8 PASS, trzy drogi dekoherencji
plus kwantowy darwinizm).
\end{theorem}

% -------------------------------------------------------------------
\subsection{Podsumowanie i~predykcje testowalne}%
\label{app:P-summary}
% -------------------------------------------------------------------

\begin{remark}[Pełny łańcuch wyprowadzenia]\label{rem:P-summary}
\begin{equation*}
  \begin{array}{c}
    \text{Aks.~\ref{ax:przestrzen}: $\Phi$ konstytuuje przestrzeń} \\
    \downarrow \\
    \text{Q0: symetria skal.\ ODE $\Rightarrow \hbar(\Phi)$} \\
    \downarrow \\
    \text{Q1: pomiar = interakcja soliton--detektor $\Rightarrow \Delta x\Delta p \geq \hbar(\Phi)$} \\
    \downarrow \\
    \text{Q2: intensywność interakcji $\sim |A|^2 \Rightarrow$ reguła Borna} \\
    \downarrow \\
    \text{Q3: linearyzacja ODE $\Rightarrow$ superpozycja} \\
    \downarrow \\
    \text{Q4: $\pi_3(S^3)$ + winding conservation $\Rightarrow$ singlet $\Rightarrow$
           $|S| = 2\sqrt{2}$} \\
    \downarrow \\
    \text{Q5--Q7: spin~$\tfrac{1}{2}$, FD/BE, dekoherencja}
  \end{array}
\end{equation*}
\textbf{Nic} nie jest postulowane poza samozwrotnością pola $\Phi$.
\end{remark}

\begin{hypothesis}[Predykcje testowalne wykraczające poza QM]%
\label{hyp:P-predictions}
TGP przewiduje odchylenia od standardowej mechaniki kwantowej
wykrywalne w~obecnych lub zaplanowanych eksperymentach:
\begin{enumerate}[nosep]
  \item \textbf{Grawitacyjna zmiana $\hbar$:}
        $\delta\hbar/\hbar \approx -GM/(rc^2) \sim 10^{-9}$
        na~Ziemi. Test: interferometria atomowa w~różnych wysokościach
        --- odchyłka widma interferencyjnego proporcjonalna
        do $\delta\hbar$.
  \item \textbf{Osłabione splątanie przy polu grawitacyjnym:}
        Korelacje Bella w~polu silniejszym dają $|S|$ nieznacznie
        mniejszy od $2\sqrt{2}$ (bo $\hbar$ mniejsze
        $\Rightarrow$ szybsza dekoherencja).
        Szacunek: $\delta|S|/|S| \sim 10^{-9}$ na~Ziemi.
  \item \textbf{Druga harmoniczna ogona:}
        $A_2/A_1 \approx 1{,}1\%$ (tw.~\ref{thm:P-born}).
        Sygnatura nieliniowości w~interferometrii precyzyjnej.
  \item \textbf{Odchylenie wykładnika Borna:}
        $p = 2{,}028$ zamiast dokładnie 2 w~reżimach silnego
        nakładania ogonów (korekcje drugiego rzędu).
\end{enumerate}
\end{hypothesis}

\begin{problem}[Program formalizacji w~Lean~4]\label{prob:P-lean}
Łańcuch Q0--Q7 pozostaje w~fazie \emph{numerycznej weryfikacji}
(79/79 PASS). Pełna formalizacja w~Lean~4 wymaga:
\begin{enumerate}[nosep]
  \item Formalnego dowodu $\pi_3(S^3) = \mathbb{Z}$
        i~$\pi_1(\mathcal{C}_{\mathrm{sol}}) = \mathbb{Z}_2$
        w~mathlib (dodatek~\ref{app:pi1-formal-proof} dostarcza szkic);
  \item Formalizacji twierdzenia spin--statystyki w~bazie substratu;
  \item Weryfikacji rachunku CHSH jako \texttt{Real.sqrt 2 * 2}
        na poziomie konstruktywnym.
\end{enumerate}
To jest program otwarty (status: O16 --- pełna formalizacja QM TGP).
\end{problem}
